\setcounter{footnote}{30}

\printedpage{29}
% head: A. Grothendieck : Dérivateurs, Ch. XVI, Édition des Universités de Montpellier II et Paris VII
% foot: Version 18 mars 2019
$\Parf_{W_0}$ respectively. (In all, this brings in six distinct sets $\mathcal{F}$, namely the six sets of diagram (4.2.1) other than $\Fib_W$ and $\Fib_{W_0}$.)

The non-tautological implication $\mathrm{C} \dotimp \mathrm{D}$ is a special case of the implication $\mathrm{C}_{W_0} \dotimp \mathrm{D}_{W_0}$, which is contained in what has just been announced, namely $\mathrm{C}_{W_0}(\mathcal{F}_{W_0}) \Longrightarrow \mathrm{C}_{W_0}(\Fib_{W_0})$, taking into account that $\Parf_{W_0} \subset \Fibf_{W_0}$ (\textit{cf.} $(*)$, p. 42). Taking
\msp{page 50}
$W_0 = W$, we thus find the implication $\mathrm{A} \Longrightarrow \mathrm{D}$, which is however contained in the other implications of the diagram, as a composite in the chain of implications
\[
\mathrm{A} \Longrightarrow \mathrm{B} \Longrightarrow \mathrm{C} \Longrightarrow \mathrm{D} \Longrightarrow \mathrm{D}_{W_0} \dotimp \mathrm{F}.
\]

It remains, in these general comments, to say a word about this last implication
\[
\mathrm{D}_{W_0} \dotimp \mathrm{F},
\]
\textit{i.e.} about the fact that the weakest of all the conditions A to $\mathrm{D}_{W_0}$ considered here already implies that Cat has enough $W$-fibrations, \textit{i.e.} is a category with cofibrations for $(W, \Fib_W)$. This is clear when we take $\mathrm{D}_{W_0}$ in the form $\mathrm{A}(\Parf_{W_0})$, \textit{i.e.}
\[
\Parf_{W_0} \subset W.
\]
As every arrow $f$ of Cat factors as $f = pi$, with
\[
i \in W_0 \subset W \quad\text{and}\quad f \in \Parf_{W_0} \subset \Fib_W\, , % [sic] f \in \Parf_{W_0} -> p \in \Parf_{W_0}, the factor of f = pi lying in Parf is p
\]
we find the factorization with $i \in W$, $f \in \Fib_W$. % [sic] f \in \Fib_W -> p \in \Fib_W, the factorisation f = pi requires p \in \Fib_W

\msp{page 51}
4.4.8. To finish establishing the implications of the diagram p. 48, it only remains to establish the equivalences
\[
\left\{
\begin{aligned}
&\mathrm{A}(\mathcal{F}_{W_0}) &&\Longleftrightarrow &&\mathrm{A}(\Fibf_{W_0})\\
&\mathrm{B}(\mathcal{F}_{W_0}) &&\Longleftrightarrow &&\mathrm{B}(\Fibf_{W_0})\\
&\mathrm{C}(\mathcal{F}_{W_0}) &&\Longleftrightarrow &&\mathrm{C}(\Fibf_{W_0}),
\end{aligned}
\right.
\tag*{(4.4.8.1)}
\]
which will give the similar equivalences for $W$, on taking $W_0 = W$.

The first equivalence (4.4.8.1) will be a consequence of the following ($\pm$ tautological) lemma:

\textbf{Lemma 4.4.9.} \textit{Let} $(\mathcal{M}, W)$ \textit{be a category of models,} $\mathcal{F}, \mathcal{G} \subset \Fl(\mathcal{M})$, \textit{such that every} $f \in \mathcal{F}$, $f : X \longrightarrow S$, \textit{can be written} $f = pi$, \textit{with} $i \in W^u_S$, $p \in \mathcal{G}$. \textit{Then} $\mathrm{A}(\mathcal{G}) \Longrightarrow \mathrm{A}(\mathcal{F})$, i.e.
\[
\mathcal{G} \subset \Fib_W \;\Longrightarrow\; \mathcal{F} \subset \Fib_W,
\tag*{(4.4.9.1)}
\]
\textit{and likewise} $\mathrm{C}(\mathcal{G}) \Longrightarrow \mathrm{C}(\mathcal{F})$, i.e.
\[
\mathcal{G} \cap W \subset W^u \;\Longrightarrow\; \mathcal{F} \cap W \subset W^u.
\tag*{(4.4.9.2)}
\]

\textsc{Proof}. A) Indeed, if $\mathcal{G} \subset \Fib_W$, the factorization $f = pi$ of $f$ is such that $p \in \Fib_W$. Now a composite $pi$, with $p \in \Fib_W$ and $i \in W^{\mathrm{u}}_S$, is in $\Fib_W$, as one verifies at once on the diagram of cartesian squares
\printedpage{30}
% head: A. Grothendieck : Dérivateurs, Ch. XVI, Édition des Universités de Montpellier II et Paris VII
% foot: Version 18 mars 2019
\[
\begin{tikzcd}
X \arrow[d, "W^{\mathrm{u}}_S"'] & X' \arrow[l] \arrow[d, "(W)"'] & X'' \arrow[l] \arrow[d, "(W)"'] \\
\overline{X} \arrow[d, "\Fib_W"'] & \overline{X'} \arrow[l] \arrow[d] & \overline{X''} \arrow[l, "(W)"'] \arrow[d] \\
S & S' \arrow[l] & S'' \arrow[l, "W"']
\end{tikzcd}
\]
on which one reads that $X'' \longrightarrow X'$ is in $W$, q.e.d.

\msp{page 52}
This lemma applies in particular to the case where $\mathcal{F} = \Fibf_{W_0}$, $\mathcal{G} = \mathcal{F}_{W_0}$, for $f$ factors as
\[
\begin{tikzcd}
X \arrow[dr, "i"] \arrow[dd, "f"'] & \\
& \Phi_S(X) \arrow[dl, "p"] \\
S, &
\end{tikzcd}
\]
and we know that $i \in W^{\mathrm{u}}_{0\,S} \subset W^{\mathrm{u}}_S$, and $p \in \mathcal{F}_{W_0} \subset \mathcal{F}_W$, which are the desired conditions.

C) Similar proof: Let $f \in \mathcal{F} \cap W$, factoring as $f = pi$, with $i \in W^{\mathrm{u}}_S$, $g \in \mathcal{G}$, we will then have $g \in W$, hence $g \in \mathcal{G} \cap W \subset W^{\mathrm{u}}$ (by hypothesis $\mathrm{C}(\mathcal{G})$), a fortiori we have $g \in W^{\mathrm{u}}_S$, hence $f \in W^{\mathrm{u}}_S$ as a composite of two arrows of $W^{\mathrm{u}}_S$, and as the target of $f$ is $S$ itself, this means $f \in W^{\mathrm{u}}$, q.e.d.
\[
\begin{tikzcd}
X \arrow[dr, "i\in W^{\mathrm{u}}_S"] \arrow[dd, "\mathcal{F}\cap W"'] & \\
& \bar{X} \arrow[dl, "g"] \\
S &
\end{tikzcd}
\]

I do not see a way of also proving
\[
\mathrm{B}(\mathcal{G}) \qquad\quad \Longrightarrow \qquad\quad \mathrm{B}(\mathcal{F}),
\]
with only the hypotheses of the preceding lemma on the pair $\mathcal{F}$, $\mathcal{G}$. It would be necessary that for every $S$ in $\mathcal{M}$, and every arrow
\msp{page 53}
$u : X \longrightarrow Y$ in $\mathcal{M}/S$, we could insert the diagram obtained in $\mathcal{M}$
\[
\begin{tikzcd}[column sep=small]
X \arrow[rr, "u"] \arrow[dr, "f"'] & & Y \arrow[dl, "g"] \\
& S &
\end{tikzcd}
\]
\printedpage{31}
% head: A. Grothendieck : Dérivateurs, Ch. XVI, Édition des Universités de Montpellier II et Paris VII
% foot: Version 18 mars 2019
into a commutative diagram below
\[
\begin{tikzcd}[column sep=small]
X \arrow[rr, "u"] \arrow[d, "i"'] & & Y \arrow[d, "j"] \\
\bar{X} \arrow[rr, "\bar{u}"] \arrow[dr, "p"'] & & \bar{Y} \arrow[dl, "q"] \\
& S &
\end{tikzcd}
\tag*{(4.4.9.3)}
\]
with
\[
f = pi, \qquad\quad g = qj, \qquad\quad i, j \in W^{\mathrm{u}}_S, \qquad\quad p, q \in \mathcal{G}.
\]
This will in any case be so if for every $S$, and a variable object $X$ of $\mathcal{M}/S$, we can find a factorization \textit{functorial in} $X$ of its structural morphism $f_X$ as $p_X i_X$, with $i_X \in W^{\mathrm{u}}_S$ and $p_X \in \mathcal{G}$, \textit{i.e.} if $i$, $p$ are obtained from a functor
\[
\varphi : \mathcal{F}_S \longrightarrow \mathcal{M}/S \ \text{\footnotemark},
\]
\footnotetext{$\mathcal{F}_S$ is the full subcategory of $\mathcal{M}/S$ consisting of the $X$ such that $f_X \in \mathcal{F}$.}
and a morphism of functors
\[
i : \underbrace{\inc_{\mathcal{F}_S}}_{\substack{\text{inclusion}\\ \text{functor}}} \longrightarrow \varphi.
\]
Assuming the existence of a diagram having the stated properties, for every arrow $u$ in $\mathcal{F}_S$, or only every arrow $u$ which is also in $W_S$, we can now prove that $u \in W_S \Longrightarrow u \in W^{\mathrm{u}}_S$, provided that the similar property is true for $\mathcal{G}$ in place of $\mathcal{F}$.

\msp{page 54}
Indeed, the hypothesis on $u \in W$ implies $\bar{u} \in W$, and the hypothesis on $\mathcal{G}$ then implies $\bar{u} \in W^{\mathrm{u}}_S$. As $W^{\mathrm{u}}_S$ is a localizer in $\mathcal{M}/S$, we do conclude $u \in W^{\mathrm{u}}_S$, q.e.d.

Thus we have obtained

\textbf{Corollary 4.4.10.} \textit{Suppose, under the conditions of lemma} 4.4.9, \textit{that every triangular commutative diagram} $\begin{tikzcd}[column sep=small, row sep=small, ampersand replacement=\&]
X \arrow[rr, "u"] \arrow[dr, "f"'] \& \& Y \arrow[dl, "g"] \\
\& S \&
\end{tikzcd}$ \textit{can be inserted into a diagram of type} (4.4.9.3) \textit{as above} (\textit{with the explicit relations on the arrows occurring in it}). \textit{Then}
\[
\mathrm{B}(\mathcal{G}) \;\Longrightarrow\; \mathrm{B}(\mathcal{F}).
\]
\textit{Thus if moreover} $\mathcal{G} \subset \mathcal{F}$, \textit{so that} $\mathrm{A}(\mathcal{F}) \Longrightarrow \mathrm{A}(\mathcal{G})$ \textit{tautologically, and likewise for} B \textit{and for} C, \textit{we find}
\[
\begin{aligned}
\mathrm{A}(\mathcal{F}) &\;\Longleftrightarrow\; \mathrm{A}(\mathcal{G})\\
\mathrm{B}(\mathcal{F}) &\;\Longleftrightarrow\; \mathrm{B}(\mathcal{G})\\
\mathrm{C}(\mathcal{F}) &\;\Longleftrightarrow\; \mathrm{C}(\mathcal{G}).
\end{aligned}
\]

We can apply this corollary in the case that concerns us here, with
\[
\mathcal{G} \overset{\text{def}}{=} \mathcal{F}_{W_0} \subset \mathcal{F} \overset{\text{def}}{=} \Fibf_{W_0},
\]
\printedpage{32}
% head: A. Grothendieck : Dérivateurs, Ch. XVI, Édition des Universités de Montpellier II et Paris VII
% foot: Version 18 mars 2019
\msp{page 55}
seeing that in this case we do have a \textit{functorial} factorization of an $f \in \Fibf_{W_0}$, $f : X \longrightarrow S$, as $f = pi$, $i \in W^{\mathrm{u}}_{0\,S} \subset W^{\mathrm{u}}_S$, $p \in \mathcal{F}_{W_0} = \mathcal{G}$, namely the one recalled on page 43:
\[
\begin{tikzcd}[column sep=small]
X \arrow[rr, "i_X"] \arrow[dr, "f=f_X"'] & & \bar{X} = \Phi_S(X) \arrow[dl, "p_X"] \\
& S. &
\end{tikzcd}
\]

This thus proves the announced equivalences (4.4.8.1) and finishes establishing the validity of the implications of the diagram on page 48.

These implications, with the sole exception of the dotted implication $\mathrm{C}_{W_0} \dottedimpl \mathrm{A}_{W_0}$ (which assumes that we have $W_0$ (7)), make no appeal to any particular hypothesis on the basic localizers $W$ and $W_0$ (apart from the inclusion $W \supset W_0$). But we can, using lemma 4.4.10, give remarkable conditional implications, under the sole condition that $W$, or $W_0$, satisfy the innocuous \textit{axiom}
\msp{page 56}
\textit{of limits} Loc (8), and even only in the weak form of stability of localizers under $\varinjlim$ of \textit{sequences}. Let us first give them in technical form, before unfolding the consequences in a theorem-scholie:

\textbf{Corollary 4.4.11.} \textit{Suppose that} $W_0$ \textit{satisfies axiom} Loc (8), \textit{or only} Loc (8 bis) (\textit{stability under} $\varinjlim$ \textit{of sequences}). \textit{Then} $\mathrm{D}_{W_0} \Longrightarrow \mathrm{A}_{W_0}$ (\textit{hence the conditions} $\mathrm{A}_{W_0}$, $\mathrm{B}_{W_0}$, $\mathrm{C}_{W_0}$, $\mathrm{D}_{W_0}$ \textit{are equivalent}). \textit{In particular} (\textit{taking} $W_0 = W$), \textit{if} $W$ \textit{satisfies} Loc (8 bis), \textit{then we have} $\mathrm{D}_W \Longrightarrow \mathrm{A}_W$ (\textit{hence the conditions} $\mathrm{A}_W$, $\mathrm{B}_W$, $\mathrm{C}_W$, $\mathrm{D}_W$ \textit{are equivalent}).

Indeed, the implication to be proved can also be written
\[
\mathrm{A}(\Parf_{W_0}) \;\Longleftrightarrow\; \mathrm{A}(\Fibf_{W_0}).
\]

The non-tautological implication $\Longrightarrow$ is covered by lemma 4.4.9. We must check that assuming Loc (8 bis) for $W_0$, every $f \in \Fib_{W_0}$, $f : X \longrightarrow S$, factors as $f = pi$, with $i \in W^{\mathrm{u}}_0$, $p \in \Parf_{W_0}$. Now we know\footnote{and this is an entirely unexpected find, a little like the lemma of K. S. \textsc{Brown} \ldots} (VI, p. 193 ff.) % [sic] f \in \Fib_{W_0} -> f \in \Fibf_{W_0}, the implication to be proved concerns A(\Fibf_{W_0}) and the next line has \Fibf_{W_0}
\msp{page 57}
that axiom Loc (8 bis) implies that $f \in \Fibf_{W_0}$ factors as
\[
\begin{tikzcd}[column sep=small]
X \arrow[rr, "i_\infty"] \arrow[dr, "f"'] & & \bar{X} = \Phi_\infty(X) \arrow[dl, "p_\infty"] \\
& S, &
\end{tikzcd}
\qquad
\text{\parbox[t]{0.33\textwidth}{\footnotesize (category of paths infinite on one side, relativized over $S$, \textit{i.e.} of pairs $(x,c)$, with $x \in \Ob X$ and $c : f(x) \longrightarrow s$ an infinite path in $S$, with source $f(x)$.)}}
\]
with $p_\infty \in \Parf_{W_0}$ (as we want), and $i_\infty \in W^{\mathrm{u}}_{0\,S}$, and as $W^{\mathrm{u}}_{0\,S} \subset W^{\mathrm{u}}_S$, we do have $i_\infty \in W^{\mathrm{u}}_S$.

\textbf{Corollary 4.4.12.} \textit{Suppose that} $W$ and $W_0$ \textit{satisfy} Loc (8 bis). \textit{Then we have} $\mathrm{D}_{W_0} \Longrightarrow \mathrm{A}$ (\textit{hence the eight conditions} A, B, C, D, $\mathrm{A}_{W_0}$, $\mathrm{B}_{W_0}$, $\mathrm{C}_{W_0}$, $\mathrm{D}_{W_0}$ \textit{are equivalent}).

\printedpage{33}
% head: A. Grothendieck : Dérivateurs, Ch. XVI, Édition des Universités de Montpellier II et Paris VII
% foot: Version 18 mars 2019
This means indeed that
\[
\mathrm{A}(\Parf_{W_0}) \;\Longrightarrow\; \mathrm{A}(\Fibf_W),
\]
and this is again covered by lemma 4.4.9. Indeed, the above factorization of an arrow $f$ does not depend on the choice of a localizer, and $p_\infty$ is $W$-perfect for \textit{every} localizer satisfying Loc (8 bis), hence in particular for $W_0$. And the arrow $i_\infty$ is in $W^{\mathrm{u}}_S$, for every localizer satisfying Loc (8 bis), hence for $W_0$.

\msp{page 58}
I am going to recapitulate the results obtained in this section in a theorem, which will constitute the main result of the present paragraph on the axioms of $W$-fibrations.

\textbf{Theorem 4.4.13.} \textit{Let} $W$ \textit{be a basic localizer in} Cat. \textit{For every localizer} $W_0$ \textit{contained in} $W$, \textit{we have made explicit a set of four conditions on the pair} $W_0$, $W$, \textit{which for} $W$ \textit{fixed, depend on} $W_0$ \textit{alone, denoted}
\[
\mathrm{A}_{W_0},\ \mathrm{B}_{W_0},\ \mathrm{C}_{W_0},\ \mathrm{D}_{W_0},
\tag*{(4.4.13.1)}
\]
\textit{and defined respectively as}
\[
\mathrm{A}(\mathcal{F}_{W_0}),\ \mathrm{B}(\mathcal{F}_{W_0}),\ \mathrm{C}(\mathcal{F}_{W_0}), \qquad\quad \textit{and}
\]
\[
\mathrm{A}(\Parf_{W_0}) \Bigl(\Longleftrightarrow\; \mathrm{B}(\Parf_{W_0}) \;\Longleftrightarrow\; \mathrm{C}(\Parf_{W_0})\Bigr),
\]
\textit{with the notations} $\mathrm{A}(\cdot)$, $\mathrm{B}(\cdot)$, $\mathrm{C}(\cdot)$ \textit{introduced in} (4.4.1). \textit{Taking} $W_0 = W$, \textit{we find the conditions on} $W$
\[
\mathrm{A},\ \mathrm{B},\ \mathrm{C},\ \mathrm{D}
\tag*{(4.4.13.2)}
\]
\textit{with} $\mathrm{A} = \mathrm{A}_W$, $\mathrm{B} = \mathrm{B}_W$, $\mathrm{C} = \mathrm{C}_W$, $\mathrm{D} = \mathrm{D}_W$, \textit{of which the first three are none other than} Loc (7 a), Loc (7 b), Loc (7 c) (cf. 4.3, \textit{page} 38).

\circled{1} \textit{For these conditions we have the following implications.}

\begin{enumerate}[leftmargin=2.5em]
\item[a)] \textit{We have the implications}
\[
\mathrm{A}_{W_0} \;\Longrightarrow\; \mathrm{B}_{W_0} \;\Longrightarrow\; \mathrm{C}_{W_0} \;\Longrightarrow\; \mathrm{D}_{W_0},
\tag*{(4.4.13.3)}
\]
(\textit{whence in particular}

\msp{page 59}
\[
\mathrm{A}_W \;\Longrightarrow\; \mathrm{B}_W \;\Longrightarrow\; \mathrm{C}_W \;\Longrightarrow\; \mathrm{D}_W\ ).
\tag*{(4.4.13.4)}
\]
\item[b)] \textit{If} $W'_0 \subset W_0$ \textit{is a basic localizer contained in} $W_0$, \textit{we have} (\textit{tautologically})
\[
\mathrm{A}_{W_0} \;\Longrightarrow\; \mathrm{A}_{W'_0}, \qquad\quad \mathrm{B}_{W_0} \;\Longrightarrow\; \mathrm{B}_{W'_0}, \qquad\quad \mathrm{C}_{W_0} \;\Longrightarrow\; \mathrm{C}_{W'_0}.
\]
\textit{We thus find, by} (a) \textit{and} (b) (\textit{taking there for the pair} $(W, W_0)$ \textit{as} $(W_0, W'_0)$),
\end{enumerate}